% Transcription — fonds Alexandre Grothendieck, shelfmark 135, batch 2 (pages 21-40).
% Established from the facsimile archives/batches/135/batch-02.pdf (a mirror of
% https://grothendieck.umontpellier.fr/135.pdf), per the skill
% transcribe-grothendieck. Pages 24, 30, 34 and 38 are blank — absent. Page 33
% is a handwritten French poem titled "Rêve", struck through with long diagonal
% strokes — not the mathematics in hand, absent. Page 36 is a 1974 calendar
% sheet from the Committee for Cultural Relations with Foreign Countries, Hanoi
% (D.R. Vietnam) — presumably how the archivists arrived at the folder's dating
% "[vers 1974]"; not mathematics, absent. Pages 39-40 open an annotated copy of
% a typescript by Hoàng Xuân Sính ("Esquisse d'une théorie des Gr-catégories",
% presented by Henri Cartan), which the inventory's "copies de tapuscrit
% annoté" announces; the typescript continues into batch 3.
% Pass: Fable 5 (claude-fable-5), 2026-08-15 — first pass, unchecked against the pages by a human.
\documentclass[11pt,a4paper]{article}
% Le préambule commun à toutes les transcriptions du fonds Grothendieck.
%
% Il tient en une idée : **l'incertitude se compose, elle ne se résout pas.**
% Une transcription qui lisse un mot douteux a détruit la seule chose pour
% laquelle on la faisait — un lecteur doit pouvoir savoir, à chaque endroit,
% ce qui était sur le feuillet et ce qui a été ajouté. Les six macros qui
% suivent sont donc l'appareil critique, et non de la décoration.
%
% Les mêmes commandes sont reconnues par `scripts/render.mjs`, qui produit la
% vue de lecture du site. Une macro ajoutée ici doit l'être aussi là-bas,
% faute de quoi le rendu échouera bruyamment — ce qui est voulu : un rendu
% silencieusement faux est pire qu'un rendu absent.

\ProvidesPackage{grothendieck}[2026/08/08 Transcriptions du fonds Grothendieck]

\usepackage{amsmath,amssymb,amsthm}
\usepackage{mathrsfs}
\usepackage{stmaryrd}
% Les diagrammes commutatifs sont partout dans ce fonds : c'est la langue
% même de la géométrie algébrique de Grothendieck, et une transcription qui
% les rendrait en prose ne transcrirait rien.
\usepackage{tikz-cd}
% Français par défaut — c'est la langue des feuillets — mais l'anglais est
% chargé aussi : la traduction et le résumé sont des documents anglais, et la
% typographie française (espaces avant les deux-points) y serait une faute.
% Ces documents basculent par \selectlanguage{english} en tête de corps.
\usepackage[english,french]{babel}
\usepackage{fontspec}
\usepackage{geometry}
\geometry{a4paper,margin=25mm,marginparwidth=28mm}
\usepackage{marginnote}
\usepackage{xcolor}
% ulem, not soul. `soul`'s \st and \ul are built on a character-by-character
% scan that XeLaTeX's font machinery breaks: the document fails with an
% undefined control sequence rather than degrading. ulem does the same two jobs
% and is Unicode-engine safe. `normalem` stops it hijacking \emph, which the
% transcriptions use for Grothendieck's own emphasis.
\usepackage[normalem]{ulem}
\usepackage{hyperref}
\hypersetup{colorlinks=true,linkcolor=black,urlcolor=black,pdfborder={0 0 0}}

% --- Métadonnées du lot -----------------------------------------------------
% Reprises telles quelles par le script de rendu, qui les affiche en tête de
% la vue de lecture. Elles fixent ce que le fichier prétend couvrir : sans
% elles, une transcription flotte sans point d'attache dans un fonds de seize
% mille feuillets.

\newcommand{\folder}[1]{\gdef\grothFolder{#1}}
\newcommand{\batch}[1]{\gdef\grothBatch{#1}}
\newcommand{\leaves}[2]{\gdef\grothFirst{#1}\gdef\grothLast{#2}}
\newcommand{\foldertitle}[1]{\gdef\grothFoldertitle{#1}}
\gdef\grothFolder{?}\gdef\grothBatch{?}\gdef\grothFirst{?}\gdef\grothLast{?}\gdef\grothFoldertitle{}

% --- Le résumé --------------------------------------------------------------
%
% Il ouvre la lecture modernisée, sous le titre « Résumé », et il est composé
% en petit. Ce n'est pas un ornement : c'est le seul passage du document qui
% ne soit pas des mathématiques, et un lecteur doit pouvoir le reconnaître
% avant de l'avoir lu — puis le sauter s'il connaît déjà le sujet, ou s'y
% arrêter s'il ne le connaît pas. Le filet qui le suit marque la reprise du
% corps.

\newenvironment{resume}
  {\small\setlength{\parskip}{0.45em}}
  {\par\medskip\hrule\medskip}

% --- Mots-clés ---------------------------------------------------------------
%
% En anglais, en clôture du résumé de la lecture modernisée : le vocabulaire
% moderne sous lequel chercher ce que les feuillets construisent. Le manifeste
% du site (`npm run manifest`) les extrait pour étiqueter la cote entière —
% cette ligne est la source unique des tags, il n'y a pas de fichier à côté.
\newcommand{\keywords}[1]{\par\smallskip\noindent
  {\footnotesize\textsc{Keywords} --- \textit{#1}}\par}

% --- L'appareil critique ----------------------------------------------------

% Le numéro de feuillet, en marge. C'est l'ancre qui permet de régler un
% désaccord sur une formule en regardant *un* feuillet plutôt qu'une chemise
% de six cents. Le site s'en sert aussi pour tourner les pages du fac-similé
% quand on fait défiler la transcription.
\newcommand{\leaf}[1]{%
  \marginnote{\footnotesize\color{gray!70}\textsf{#1}}%
  \label{leaf:#1}\ignorespaces}

% Illisible : on ne devine pas. Un mot inventé qui se lit comme les autres est
% le pire résultat possible.
\newcommand{\ill}{\textcolor{red!60!black}{[\,\ldots\,]}}

% Lecture douteuse : le mot est donné, et signalé comme incertain.
\newcommand{\uncertain}[1]{\uline{#1}}

% Ajout de l'éditeur — une lettre manquante, un mot sous-entendu.
\newcommand{\add}[1]{\textcolor{blue!50!black}{[#1]}}

% Biffé par Grothendieck lui-même. Ce qu'il a rayé fait partie du document :
% une rature dit souvent où la pensée a changé de direction.
\newcommand{\struck}[1]{\sout{#1}}

% Note du transcripteur, sur l'état matériel du feuillet.
\newcommand{\note}[1]{\par\smallskip{\small\color{gray!60!black}\textit{[#1]}}\par\smallskip}

% Note marginale de Grothendieck, distincte du corps du texte.
\newcommand{\marginal}[1]{\par\smallskip\noindent\hspace{1em}%
  {\small\color{gray!60!black}#1}\par\smallskip}

% --- Titre ------------------------------------------------------------------

\AtBeginDocument{%
  \begin{center}
    {\large\bfseries Fonds Alexandre Grothendieck --- cote n\textsuperscript{o}~\grothFolder}\\[2pt]
    {\itshape\grothFoldertitle}\\[4pt]
    {\small feuillets \grothFirst--\grothLast\ (lot \grothBatch)}\\[2pt]
    {\footnotesize\color{gray!60!black}Transcription établie d'après le fac-similé numérisé
     par l'Université de Montpellier.\\
     Les lectures incertaines sont soulignées, les illisibles notées [\,\ldots\,],
     les ajouts entre crochets.}
  \end{center}
  \vspace{1em}}

\endinput


\folder{135}
\batch{2}
\pages{21}{40}
\dating{[vers 1974]}
\watermark{Édition de démonstration}
\foldertitle{Catégories [et Gr-catégories] : notes manuscrites (s.d.), copies de tapuscrit annoté (s.d.), copies de manuscrit annoté (s.d.).}

\begin{document}

\section*{Extensions d'une catégorie par un groupe, suite (pages 21 à 23)}

\page{21}
\note{feuille sans pagination de l'auteur, couverte de calculs que séparent de
longues diagonales ; elle prépare visiblement la note des pages 20, 22 et 23 ;
nous transcrivons zone par zone, de haut en bas}

\[
P_u \otimes P_v \otimes P_w \longrightarrow P_{uvw}
\]

\uncertain{$1 \times Q \otimes R \to S$}, puis les quatre formes de
l'accouplement :
\[
P \otimes Q \to S \otimes R^{-1}, \qquad
Q \otimes R \to P^{-1} \otimes S,
\]
\[
R \otimes S^{-1} \to Q^{-1} \otimes P^{-1}, \qquad
S^{-1} \otimes P \to R^{-1} \otimes Q^{-1}.
\]
\note{les quatre lignes sont réunies par une accolade ; les exposants $^{-1}$
sont par endroits surchargés et restent douteux}

La colonne simpliciale, chaque étage envoyé sur le suivant par un faisceau de
barres verticales (les flèches de faces) :
\[
\mathrm{Fl}_3\,F, \qquad
M_{\Gamma_2} \to \mathrm{Fl}_2(\Gamma) = \mathrm{Fl}\,\Gamma
\times_{\mathrm{Ob}} \mathrm{Fl}\,\Gamma, \qquad
M_{\Gamma_1} \to \mathrm{Fl}_1(\Gamma) = \mathrm{Fl}\,F, \qquad
M_{\Gamma_0} \to \mathrm{Fl}_0(\Gamma) = \mathrm{Ob}\,\Gamma
\]
\note{l'auteur passe de $F$ à $\Gamma$ d'une ligne à l'autre ; le premier
indice de la deuxième ligne, surchargé, se lit aussi bien $M_{P_2}$ ; la base
du produit fibré n'est pas écrite en toutes lettres}
\[
H^2(\Gamma, M), \qquad H^1(\Gamma, M), \qquad H^0(\Gamma, M),
\]
et, isolé plus bas, un $H^2$.

L'échelle parallèle, à droite, avec les mêmes barres de faces entre étages —
en marge des deux premiers étages : $L_u \otimes L_v \otimes L_w \to L_{uvw}$
et $L_u \otimes L_v \to L_{uv}$ :

\begin{itemize}
\item $\Gamma_3 \to P_3$ = \uncertain{accouplements de trois objets vers un
quatrième} ;
\item $\Gamma_2 \to P_2$ = \uncertain{accouplements de deux objets de $P$ vers
un troisième} ;
\item $\Gamma_1 \to P_1 = \mathrm{Ob}\,P$ ;
\item $\Gamma_0 \to e$.
\end{itemize}

$P_I$ = ensemble des \struck{\ill} \uncertain{triples} $(L, M, \varphi)$ :
$L : I \to \mathrm{Ob}\,P$, soit $L = (L_i)_{i \in I}$ ; $M \in \mathrm{Ob}\,P$ ;
$\varphi : \bigotimes_{i \in I} L_i \to M$.
\note{sous le $P_I$, un renvoi de deux mots que nous ne savons pas lire ;
à droite, \uncertain{$\Gamma_I \to \Pi_1(P)$}}

$L_u L_v \simeq L_w$ — le petit triangle en marge fixe les notations : deux
flèches composables $u$, $v$, la composée $w$, le côté $u$ portant $L_u$.
\[
L_{.} \otimes L_{.} \longrightarrow M, \qquad
L_{.} \otimes L_{.} \otimes M^{-1} \simeq\ \ill{}, \qquad
L_{.} \otimes \cdots \otimes L_{.} \otimes S^{-1} \simeq 1.
\]
\note{les indices des $L$ de ces trois lignes sont réduits à des points ou
surchargés}

Les trois formes d'un accouplement :
\[
P \otimes Q \to R, \qquad
R \otimes Q^{-1} \to P, \qquad
P^{-1} \otimes R \to Q.
\]
\note{à droite, deux lignes griffonnées — on y lit un $Q \otimes R$ — que nous
ne savons pas raccorder}

En bas à gauche, un tétraèdre esquissé : quatre sommets, les six arêtes
portant $u$, $v$, $w$, $u'$, $v'$, $w'$, et le long de deux côtés $L_u$ et
$L_v$. \note{l'orientation de chaque arête ne se laisse pas fixer avec
certitude ; nous ne dessinons pas ce que nous ne savons pas lire}

Enfin les relations de cohérence, sur deux colonnes :
\[
L_u L_v \sim L_{uv}, \qquad
L_{uv} L_w \sim L_{uvw}, \qquad
L_u L_v L_w \sim L_{uvw} ;
\]
\[
L_v L_w \sim L_{vw}, \qquad
L_u L_{vw} \sim L_{uvw}, \qquad
L_u L_v L_w \sim L_{uvw}.
\]
\note{il s'agit visiblement des deux parenthésages du produit triple ; les
indices exacts, serrés et par endroits surchargés, restent douteux}

\page{22}
\note{page 2 de la main de l'auteur : la note commencée page 20 se poursuit,
l'axiome c) des 2-extensions s'achevant sur son diagramme}

\begin{tikzcd}[column sep=small, row sep=small, nodes={font=\scriptsize}]
(P_u \otimes P_v) \otimes P_w \arrow[r, "\varphi_{u,v} \otimes \mathrm{id}"]
\arrow[d, "\mathrm{can}"'] & P_{uv} \otimes P_w \arrow[dd, "\varphi_{uv,w}"] \\
P_u \otimes (P_v \otimes P_w) \arrow[d, "\mathrm{id} \otimes \varphi_{v,w}"'] & \\
P_u \otimes P_{vw} \arrow[r, "\varphi_{u,vw}"'] & P_{uvw}
\end{tikzcd}

\note{entre les deux composés, la 2-flèche $\alpha_{u,v,w}$}

d) Un axiome de compatibilité \uncertain{pour} \struck{\ill} quatre flèches
composables. On veut en fait que tous les isomorphismes entre
\uncertain{foncteurs} de la forme
\[
\bigl(\cdots((P_{u_1} \otimes P_{u_2}) \otimes P_{u_3}) \cdots\bigr)
\longrightarrow P_{u_1 u_2 \cdots u_n},
\]
déduits des $\varphi_{u,v}$ \uncertain{et des} isom.\ d'associativité,
\uncertain{forment un système transitif}.
\note{la phrase est chargée de ratures ; en marge gauche, une colonne biffée
d'essais de parenthésages — $((P_t \otimes P_u) \otimes P_v) \otimes P_w$,
puis $(P_t \otimes (P_u \otimes P_v)) \otimes P_w$, chaque flèche marquée
« can » — et un long trait ondulé menant à $P_{u_1 u_2 \cdots u_n}$ et à
$P_t \otimes (P_u \otimes P_v \otimes P_w)$}

NB — Au lieu de $M$, on peut aussi se donner une 2-\struck{Gr-}cat.\
\uncertain{de Picard} $\widetilde{P}$ (p.ex.\ celle des $M$-gerbes)] :
\[
\mathrm{Fl}_1(\Gamma) \longrightarrow \pi_0(\widetilde{P}), \qquad
\mathrm{Fl}_2(\Gamma) \longrightarrow \ill{}
\]
\note{un grand crochet droit court le long du NB et des deux formules ; la
parenthèse qui suit $\pi_0(\widetilde{P})$ ne se laisse pas lire}

\page{23}
\note{page 3 de la main de l'auteur ; les numéros 1) à 3) sont cerclés sur la
page}

$\Gamma$ cat., $P$ Gr-cat.
\[
\mathrm{Fl}(\Gamma) \xrightarrow{\ \alpha\ } \pi_0(P) \longrightarrow
\mathrm{Aut}(\pi_1(P)),
\]
la composée étant notée $\overline{\alpha}$. On va définir la 2-catégorie des
$\alpha$-extensions de $\Gamma$ par $P$ :

1) 0-objets : les ext.\ de $\Gamma$ par $P$ subordonnées à $\alpha$ ;

2) 1-objets : $\{P_u, \varphi_{u,v}\} \to \{P'_u, \varphi'_{u,v}\}$ sont les
systèmes d'hom
\[
\lambda(u) : P_u \longrightarrow P'_u
\]
donnant lieu à diagrammes commutatifs \struck{avec les \ill{} $\Pi_{u,v,w}$ et
$\Pi'_{u,v,w}$} :

\begin{tikzcd}[column sep=normal, row sep=normal]
P_u \otimes P_v \arrow[r, "\varphi_{u,v}"]
\arrow[d, "\lambda(u) \otimes \lambda(v)"'] & P_{uv} \arrow[d, "\lambda(uv)"] \\
P'_u \otimes P'_v \arrow[r, "\varphi'_{u,v}"'] & P'_{uv}
\end{tikzcd}

NB — Si on a un tel $\{\lambda(u)\}$, les autres \struck{$\mu'(u)$}
$\{\mu(u)\}$ sont de la forme $\lambda(u) \otimes f(u)$,
$f : \mathrm{Fl}\,\Gamma \to \pi_1(P)$, avec
\[
f(uv) = f(u).(\alpha(u).f(v)),
\]
i.e.\ $f$ définit une 1-cochaîne sur $\Gamma$ à valeurs dans $(\pi_1(P),
\overline{\alpha})$. Ainsi l'ens.\ des hom de $\{P_u, \varphi_{u,v}\}$ en
$\{P'_u, \varphi'_{u,v}\}$ est un \struck{\ill} pseudo-torseur sous
$Z^1(\Gamma, (\pi_1(P), \overline{\alpha}))$.

3) 2-objets : entre $(\lambda_u), (\mu_u) : \{P_u, \varphi_{u,v}\}
\rightrightarrows \{P'_u, \varphi'_{u,v}\}$. Donc $\mu_u = \lambda_u \otimes
f(u)$, $f(u) \in Z^1(\Gamma, \pi_1(P))$. On \uncertain{pose} que l'ens.\ des
flèches de $\lambda$ en $\mu$ est l'ens.\ des $m \in Z^0(\Gamma, \pi_1(P))$
tels que $f = \delta m$, i.e.
\[
\mu_u = \lambda_u \otimes (m(y) - \alpha(u).m(x)) \qquad (?)
\]
\note{le point d'interrogation final est de l'auteur ; le sous-indice du
$\mu$ de gauche, surchargé, n'est plus lisible ; en marge, le croquis
$x \xrightarrow{\ u\ } y$ fixe les notations}

\section*{Un programme : $n$-catégories et $n$-champs (pages 25 à 29)}

\page{25}
\note{les pages 25 à 29 forment une seule liste numérotée de 1 à 11 — un
programme d'étude, d'une écriture serrée et très abrégée}

1. $n$-catégories (strictes — pas strictes) et $n$-foncteurs (stricts ou
pas). \uncertain{Il y a des chances que les « pas strictes » se ramènent aux
strictes.} Les $n$-catégories forment une $(n{+}1)$-catégorie, les Hom y sont
des $n$-catégories ; notion de $i$-fidélité ($0 \leq i \leq n+1$) — la
$(n{+}1)$-fidélité \uncertain{s'appelle} $n$-équivalence.
\uncertain{Relations entre $n$-catégories et \ill{}, catégories réduites
(dans les 2 sens).} Une Gr-catégorie « est » une 2-catégorie ayant un seul
0-objet.

2. Exemples liés aux structures semi-simpliciales.

Espaces topologiques, groupes top., anneaux top.\ etc. Ensembles
($\tfrac12$-)simpliciaux \struck{\ill} (\uncertain{de Kan} ?), groupes,
\uncertain{anneaux, en connexion étroite} \uncertain{itérés}.
\uncertain{Abstractisation} : catégorie $\tfrac12$-simpliciale.

$A$ catégorie, $C$ \struck{\ill} catégorie (p.ex.\ ens.\ finis tot.\
ordonnés \uncertain{non vides} ...) : la donnée d'un foncteur
$\mathrm{Hom}. : A^{\circ} \times A \to \hat{C}$
\marginal{$(A,B) \mapsto \mathrm{Hom}.(A,B)$}
équivaut à celle d'un foncteur $A^{\circ} \times A \times C^{\circ} \to$
\uncertain{Ens}, ou encore $C^{\circ} \times A \xrightarrow{\mathrm{Hom}}
\hat{A}$, $(\sigma, B) \mapsto \mathrm{Hom}(\sigma, B)$, ou d'un foncteur
\marginal{$\sigma \mapsto \varphi_{\sigma}$}
$C^{\circ} \to \mathrm{Hom}(A, \hat{A})\ [\cong \mathrm{Hom}_{!}(\hat{A},
\hat{A})]$ (le $!$ indique les foncteurs qui commutent aux $\varinjlim$
\uncertain{quelc.}), soit encore un foncteur $\hat{C}^{\circ} \to
\mathrm{Hom}(A, \hat{A}) \simeq \mathrm{Hom}_{!}(\hat{A}, \hat{A})$
commutant aux $\varinjlim$ quelc., qu'on note encore \uncertain{$\sigma
\mapsto h_{\sigma}$}.

Si $C$ a objet final $e$ (\uncertain{exemple dans le cas précisé}) alors la
donnée d'un isom bifonctoriel
\[
(\alpha) \qquad \mathrm{Hom}.(A,B)(e) \simeq \mathrm{Hom}(A,B)
\]
équivaut à la donnée d'un isom.
\[
(\beta) \qquad \varphi_{e} = \mathrm{id}_{\hat{A}}.
\]

La donnée d'accouplements de bifoncteurs
\[
(a) \qquad \mathrm{Hom}.(A,B) \times \mathrm{Hom}.(B,C) \longrightarrow
\mathrm{Hom}.(A,C)
\]
équivaut sauf erreur à celle d'\uncertain{homomorphismes} bifonctoriels
\[
(b) \qquad \varphi_{K \times L} \Longleftarrow \varphi_{K} \circ \varphi_{L}
\]
(ou encore, d'hom fonctoriels
\uncertain{$(k)\ \varphi_{K} \longleftarrow \varphi_{M} \circ \varphi_{K}$}),
et la \uncertain{condition} d'un \uncertain{isom} d'associativité pour les
accouplements entre les Hom.\ équivaut sauf erreur à celle d'une
\struck{\ill}compatibilité d'associativité pour les données \uncertain{
ci-dessus}. Compatibilités
\marginal{$(b')$ quand \uncertain{$\varphi$} est-il un isomorphisme ?
$\mathrm{Hom}.(A, \mathrm{Hom}(L,B)) \to \mathrm{Hom}(L, \mathrm{Hom}.(A,B))$
\uncertain{est un isom.}}

\page{26}
entre données $\alpha$ et $a)$, resp.\ $\beta$ et $b)$ (\uncertain{ou}
$c)$).

Dans le cas $C$ = ens.\ tot.\ ordonnés $\neq \emptyset$, i.e.\ $\hat{C}$ =
ens.\ $\tfrac12$-simpl., on a dans les $\mathrm{Hom}.(A,B)$ des
\uncertain{notions} d'homotopie. Cela permet d'associer à $A$ des
$n$-catégories pour tout $n$ (construction facile si les $\mathrm{Hom}.(A,B)$
sont des ens.\ $\tfrac12$-simpliciaux de Kan ...).

\struck{Hom} Définition de $\mathrm{Isom}.(A,B)$ et $\mathrm{Aut}.(A)$ (le
\uncertain{premier} un \uncertain{groupe} de $\hat{C}$, le
\uncertain{deuxième} un pseudo-torseur sous celui-ci).
\note{l'ordre attendu est l'inverse — le groupe serait $\mathrm{Aut}.(A)$ et
le pseudo-torseur $\mathrm{Isom}.(A,B)$ ; la page porte ce qu'elle porte}
Notions d'homotopie de morphismes, d'isotopie, d'isomorphismes. Foncteurs
$\tfrac12$-simpliciaux des catégories à structure
\struck{semi-simpliciales} (groupoïdes à structure semi-simpliciale
etc....). \note{« à structure » est un ajout interlinéaire} On peut dire
quand un tel foncteur est une [$n$-]équivalence d'homotopie (s'il induit sur
les Hom.\ des [$n$-]équivalences d'homotopie) ; sauf erreur, cela signifie
aussi que le $n$-foncteur entre les $n$-catégories associées est une
$n$-équivalence. \uncertain{Dold-Puppe.} \uncertain{Cas des complexes dans
cat.\ abéliennes :} \uncertain{lien avec ce qui précède.}
\note{ces derniers mots, en interligne serré, restent pour partie
conjecturaux}

2'. Topos de Quillen \uncertain{et interprétation} des $n$-catégories
strictes comme structures \uncertain{$\tfrac12$-simpliciales particulières}.

3. $n$-catégories fibrées sur une catégorie, $n$-champs sur un site ou
topos. \uncertain{Ils forment une $(n{+}1)$-catégorie.} Il faudrait voir :

\begin{enumerate}
\item il y a une \uncertain{$(n{+}1)$-}équivalence entre $n$-champs sur un
site, et $n$-champs sur le topos associé ;
\item toute $n$-catégorie fibrée sur un site définit un $n$-champs
enveloppant ;
\item existence d'images inverses de $n$-champs (par morphismes (de sites
ou) de topos). [NB — La notion d'image directe « est claire » ; celle
d'images inverses \uncertain{s'en déduit} par \uncertain{formule d'}adjonction
\[
\mathrm{Hom}(A, f_{*}(B)) \simeq \mathrm{Hom}(f^{*}(A), B)
\]
{[}équivalence de $n$-catégories....].
\end{enumerate}

\page{27}
4. \struck{\ill} $n$-champs de \uncertain{profondeur} \uncertain{sur $Y$}
$\geq i$ ($0 \leq i \leq n+2$), $Y$ étant un fermé de $X$.

5. Formules de \uncertain{chgt de base} d'un top.\ ordinaire à
\uncertain{ou} une topologie étale... (Il y en a au moins trois
\struck{\ill} types \uncertain{principaux}....)
\marginal{$n$-champs ind-finis constructibles etc....}

6. \struck{\ill} Théorèmes du type de Lefschetz en topologie étale
(schémas) ou dans l'analytique-complexe.

7. $n$-champs localement constants.

Un morphisme de topos $f : X \to Y$ est une équivalence d'homotopie ssi pour
tout $n$-champ loc.\ ct.\ $A$ sur $Y$, il induit une équivalence
\[
\Gamma(Y, A) \simeq \Gamma(X, f^{*}(A)).
\]
\note{le signe $\simeq$ porte un exposant entre parenthèses, surchargé —
$(n)$ ?}
On peut donner une version \uncertain{avec} les types d'homotopie
« complétés » \uncertain{pour une catégorie de groupes} : le Artin-Mazur ?

[Cela suggère qu'un « type d'homotopie » de topos est \uncertain{reconnu}
par la connaissance, pour tout $n$, de la \uncertain{$(n{+}1)$}-catégorie des
$n$-champs loc.\ constants dessus, avec les foncteurs de passage des
\struck{\ill}champs aux $(n{-}1)$-champs sous-jacents et
\uncertain{inversement}. Dans le cadre d'Artin-Mazur, il faudrait des
restrictions de constructibilité convenables sur les $n$-champs avec
lesquels on travaille...]

8. Structures algébriques sur des $n$-catégories et $n$-champs :
Gr-structures, structure de Picard, de Picard-stricte, torseurs etc.
Stabilité \uncertain{par} images directes et inverses de \uncertain{champs}.

9. La \uncertain{$(n{+}1)$}-catégorie des $n$-champs de Picard stricts sur
$X$ est \uncertain{$(n{+}1)$}-équivalente à la sous-catégorie pleine de

\page{28}
la catégorie dérivée $D(\mathrm{Ab}(X))$ formée des complexes \struck{qui
sont} tels que $\underline{H}^{i}(K^{\cdot}) = 0$ si $i \notin [0,n]$. Les
opérations, images inverses et images directes \struck{\ill} sur les
$n$-champs de Picard stricts correspondant aux opérations analogues dans les
catégories dérivées (sauf qu'il faut tronquer, dans le cas de l'image
directe). Les Hom de cat.\ de Picard strictes \uncertain{sur $[0,n]$}
correspondent aux $R\mathrm{Hom}$ tronqués dans les catégories dérivées.

NB — Il y aurait lieu d'interpréter de façon géométrique les objets
tronqués de $D(\mathrm{Mod}(O))$, $O$ un Anneau sur $X$, en termes de cat.\
de Picard strictes \uncertain{où $O$ opère}, et d'interpréter
géométriquement les opérations $Rf_{*}$, et $Lf^{*}$ et $R\mathrm{Hom}$.
Interprétation des $\otimes^{(L)}$ en termes d'accouplements
\uncertain{« universels »} des cat.\ de Picard ?
\note{le $\otimes$ porte le $(L)$ en exposant, et au-dessous un indice
biffé}

10. Étude des $n$-Gr-champs. J'ai regardé le cas $n = 1$, $X$ quelconque
(\uncertain{cf lettre à Deligne}). Mais si $X$ est le topos ponctuel je ne
sais quoi dire pour le cas $n = 2$ p.ex. Il faudrait examiner le cas de la
$n$-catégorie associée \uncertain{à} : un groupe $\tfrac12$-simplicial
(p.ex.\ via un groupe topologique), cf.\ §2, \uncertain{étant} la
$n$-catégorie des \uncertain{homomorphismes} de \struck{\ill}
\uncertain{l'objet $\tfrac12$-simplicial unité} dans ledit groupe
$\tfrac12$-simplicial (pas strict !).

11. $n$-\struck{catégorie}\uncertain{champ} de Picard \uncertain{enveloppant}
d'un \uncertain{$(n)$-$\otimes$-champ} ACU. Ses invariants $\pi_{i}$ sont
les faisceaux $K_{i}(A)$.

\page{29}
\textbf{Question} — Que peut-on dire de la structure des $n$-champs de
Picard pas stricts ? On voit pour $n = 1$ (cf.\ correspondance avec Breen et
Deligne), mais j'ignore pour $n \geq 2$ ....

\section*{Tours de lacets et suites de fibrations, I (page 31)}

\page{31}
\note{feuille de calcul sans un mot de prose : des treillis de suites de
fibrations, à sept lettres — $S$, $Y$, $G$, $F$, $X$, $Z$, $e$ — que nous
transcrivons par blocs. Le motif de quatre rangées ci-dessous se répète sur
la page une dizaine de colonnes durant (les exposants de $\Omega$ croissant
vers la droite jusqu'à $\Omega^{5}$), puis une seconde fois en entier, puis
une troisième fois sur trois rangées s'achevant sur une colonne $e$ ; un
cadre tracé à main levée isole au centre un exemplaire du motif ; quelques
croisements portent un carré plein, des doubles barres identifient
verticalement certaines cases voisines, et un pâté d'encre biffe un essai
sous le deuxième bloc}

\begin{tikzcd}[column sep=small, row sep=small, nodes={font=\scriptsize}]
G \arrow[d] & \Omega S \arrow[l] \arrow[d] & \Omega Y \arrow[l] \arrow[d] &
\Omega G \arrow[l] \arrow[d] & \Omega^{2} S \arrow[l] \arrow[d] &
\Omega^{2} Y \arrow[l] \arrow[d] & \Omega^{2} G \arrow[l] \arrow[d] \\
e \arrow[d] & F \arrow[l] \arrow[d] & F \arrow[l, Rightarrow] \arrow[d] &
e \arrow[l] \arrow[d] & \Omega F \arrow[l] \arrow[d] &
\Omega F \arrow[l, Rightarrow] \arrow[d] & e \arrow[l] \arrow[d] \\
X \arrow[d] & Z \arrow[l] \arrow[d] & G \arrow[l] \arrow[d] &
\Omega X \arrow[l] \arrow[d] & \Omega Z \arrow[l] \arrow[d] &
\Omega G \arrow[l] \arrow[d] & \Omega^{2} X \arrow[l] \arrow[d] \\
S & Y \arrow[l] & G \arrow[l] & \Omega S \arrow[l] & \Omega Y \arrow[l] &
\Omega G \arrow[l] & \Omega^{2} S \arrow[l]
\end{tikzcd}

\note{en fin de rangées, plusieurs exposants surchargés restent douteux}

Au bas de la feuille, des essais autour d'un carré $X_{ij}$ :
\[
X_{00}, \qquad
X_{20} \times X_{02} \longleftarrow Q \longleftarrow X_{22},
\]
\note{le $X_{02}$ du produit est surchargé d'une rature}
\[
Q \longleftarrow P \longleftarrow P \times \Omega Q \longleftarrow
\Omega Q \longleftarrow \Omega P \longleftarrow \Omega P \times \Omega^{2} Q,
\]

\begin{tikzcd}[column sep=small, row sep=small, nodes={font=\scriptsize}]
e \arrow[d] & X_{12} \arrow[l] \arrow[d] & \\
e \arrow[d] & X_{10} \arrow[l] \arrow[d] & \bullet \arrow[l] \arrow[d] \\
X_{00} & X_{10} \arrow[l] & X_{11} \arrow[l]
\end{tikzcd}

\note{la case marquée $\bullet$ est entièrement surchargée — un $X_{12}$
récrit sur autre chose ? — et le $X_{10}$ de la dernière rangée se lit aussi
bien $X_{01}$}
\[
X_{20} \times X_{11} \times X_{02} \longleftarrow Q \longrightarrow\ ?,
\qquad
X_{00} \longleftarrow X_{11} \longleftarrow Q.
\]
\note{le point d'interrogation est de l'auteur}

Enfin trois lignes reliées verticalement par des doubles barres sous les
$Q$ :
\[
X_{20} \longleftarrow Q \longleftarrow \uncertain{X_{12}} \longleftarrow
\Omega X_{20},
\]
\[
X_{20} \times X_{02} \longleftarrow Q \longleftarrow X_{22} \longleftarrow
\uncertain{\Omega X_{10} \times \Omega X_{02}},
\]
\[
X_{02} \longleftarrow Q \longleftarrow X_{21} \longleftarrow \ill{}
\]
\note{la fin de la dernière ligne est griffonnée puis biffée}

\section*{Catégorie de Picard-type non stricte (page 32)}

\page{32}
\textbf{Catégorie de Picard-type non stricte.}

$A$ anneau commutatif (avec 1), \struck{\ill} $2 \cdot 1_{A} \neq 0$.

$C$ : $\mathrm{Ob}\,C$ = Modules gradués \uncertain{mod 2} sur $A$, munis
d'une \struck{base} \uncertain{formée} d'un élément, mais
\uncertain{donnée seulement au signe près} \uncertain{(soit : deux éléments,
échangés par $x \mapsto -x$, formant des bases)} ;
\note{« mod 2 » est écrit en interligne ; la parenthèse, très abrégée sur la
page, reste en partie conjecturale}

$\mathrm{Hom}(L,M)$ = ensemble des \uncertain{homomorphismes}
\uncertain{gradués} $\varphi : L \to M$ tels que $\varphi(e_{L}) = e_{M}$ ;

$\otimes$ = produit tensoriel gradué, avec
\[
e_{(L \otimes M)} = \{\, x \otimes y \mid x \in e_{L},\ y \in e_{M} \,\} ;
\]

donnée d'associativité : comme d'habitude \uncertain{(compatible d'après la
structure des bases)} ;

donnée de commutativité : règle de Koszul.

\textbf{Description alternative} : \struck{\ill} $C = (\mathbf{Z}/2\mathbf{Z})
\times C_{0}$, le premier facteur catégorie discrète, le second la catégorie
des torseurs sous $\mathbf{Z}/2\mathbf{Z}$.
\note{une accolade réunit le produit sous la légende « torseurs gradués
mod 2 sous $\mathbf{Z}/2\mathbf{Z}$ »}

\note{tout le bas de la page est barré de longues diagonales ; l'intitulé
« Torseurs gradués », souligné, est lui-même rayé ; nous transcrivons le
passage tel quel}

\struck{Torseurs gradués. — $F$ faisceau abélien sur topos $S$ (écrit
additivement) ; $\Gamma$ groupe discret ; $\Gamma \times \Gamma \to
H^{0}(S,F)$ application \ill{}.}

\struck{$C = H^{0}(S, \Gamma) \times \mathrm{Tors}(F)$ — le premier facteur
catégorie discrète d'ens.\ sous-jacent $H^{0}(S, \Gamma)$, le second la
catégorie des torseurs sous $F$ ; $C \times C \xrightarrow{\ \otimes\ } C$,
$(\alpha, P) \otimes (\beta, Q) = (\alpha\beta,\ P \otimes Q)$ ;
associativité :
$((\alpha,P) \otimes (\beta,Q)) \otimes (\gamma,R) \simeq (\alpha,P) \otimes
((\beta,Q) \otimes (\gamma,R))$, soit $(\alpha\beta\gamma,\ (P \otimes Q)
\otimes R)$ et $(\alpha\beta\gamma,\ P \otimes (Q \otimes R))$ — provient de
l'associativité de $\otimes$ dans $\mathrm{Tors}(F)$.}
\note{l'argument du second $H^{0}$ est abrégé sur la page — $\Gamma_{S}$ ? —
et plusieurs lettres grecques de la ligne d'associativité sont surchargées}

\section*{Tours de lacets et suites de fibrations, II (pages 35 et 37)}

\page{35}
\note{grande feuille pliée, couverte de diagrammes. Nous transcrivons dans
l'ordre : l'escalier de suites de fibrations attaché à des flèches
composables, les conditions équivalentes notées en clair au centre, les deux
remarques cerclées du bas, et les formules sur les Hom de foncteurs ; une
vingtaine de petits carrés d'essai, trop surchargés, restent hors
d'atteinte}

L'escalier, pour des flèches composables $A \xrightarrow{f} B
\xrightarrow{g} C$ (puis $h$, avec un quatrième objet $D$) et leurs fibres
$\Gamma$ — chaque suite se déduisant de sa voisine en appliquant $\Omega$,
et l'escalier descendant jusqu'aux flèches elles-mêmes ; en voici les
marches lisibles :
\[
\Omega\Gamma_{f} \to \Omega\Gamma_{gf} \to \Omega A \to e, \qquad
\Omega\Gamma_{g} \to \Omega D \to \Gamma(f) \to e,
\]
\[
\Omega C \to \Gamma(gf) \to \Gamma(g) \to e, \qquad
\Omega D \to \Gamma(hgf) \to \Gamma(hg) \to \Gamma(h) \to e,
\]
\note{plusieurs indices composés — $gf$, $hg$, $hgf$ — sont serrés au point
d'être conjecturaux ; l'escalier se prolonge vers le haut en $\Omega^{2}$ et
$\Omega^{3}$}
et, symétriquement, les mêmes suites en $\Sigma$ :
\[
e \leftarrow \Sigma\Gamma_{f} \leftarrow \Sigma\Gamma_{gf} \leftarrow
\Sigma\Gamma_{hgf} \to \Sigma A, \qquad
\Sigma\Gamma_{g} \to \Sigma\Gamma_{hg} \to \Sigma B \to
\Sigma^{2}\Gamma_{f},
\]
\[
\Sigma\Gamma_{h} \to \Sigma C \to \Sigma^{2}\Gamma_{gf}, \qquad \Sigma D.
\]

Au centre gauche, le carré d'un produit fibré $Z$ au-dessus de $S$, avec ses
projections $p$, $q$, la flèche $r$ et les fibres $\Gamma_{f}$,
$\Gamma_{g}$, $\Gamma_{p}$, prolongé en $\Omega X$, $\Omega S$ ; puis, en
clair :

\uncertain{Conditions équivalentes à « $S$ connexe »} (l'équiv.\ de a) et
b) \ill{}) :

\begin{enumerate}
\item $\Gamma_{g} \to \Gamma_{f}$ \uncertain{isom} ;
\item $\Gamma_{p} \to \Gamma_{g}$ \uncertain{isom} ;
\item $\forall\, T$, \uncertain{$p' : T \to X$, $q' : T \to Y$} avec
$f p' = g q'$, $\exists\, u : T \to Z$ avec $p' = pu$, $q' = qu$ (NB — $u$
n'est pas \uncertain{nécessairement} unique \ill{}) ;
\item $\Gamma_{r} \rightrightarrows \Gamma_{f}, \Gamma_{g}$
\uncertain{font de $\Gamma_{r}$ un produit fibré de $\Gamma_{f}$ et
$\Gamma_{g}$}.
\end{enumerate}
\note{les étiquettes de la liste sont sur la page a), b), [a)], c) — le
troisième item est encadré de crochets ; la fin du NB, trois lignes raturées
et récrites, ne se laisse pas suivre}

En bas, deux remarques cerclées :
\begin{itemize}
\item \uncertain{Données de $\Omega$, et des notions de suite distinguée}
\[
\Omega S \to F \to X \to S
\]
\uncertain{avec axiomes (y compris \ill{})} ;
\item \uncertain{Axiomatique des \ill{}} \struck{\ill}.
\end{itemize}

Et les formules sur les Hom de foncteurs :
$D \xrightarrow{\ \varphi\ } \hat{C}$ \uncertain{avec lim finies} \ill{},
$\mathrm{Hom}\uncertain{{}_{sex}}(D, \hat{C})$ :
\[
\mathrm{Hom}(\varphi, \psi) \in \hat{C}, \qquad
\mathrm{Hom}(\varphi, \psi)(S) =
\mathrm{Hom}(\uncertain{i_{S} \circ \varphi},\ \uncertain{i_{S} \circ \psi}).
\]
NB — Si $C$ \uncertain{est une $U$-cat}, et si \uncertain{$\varphi$, $\psi$
prennent leurs valeurs dans $\hat{C}$}, alors $\mathrm{Hom}(\varphi, \psi)
\in \hat{C}$.

NB — $\mathrm{Hom}(E, \mathrm{Hom}(\varphi, \psi)) =
\mathrm{Hom}(\uncertain{i_{E} \circ \varphi},\ \uncertain{i_{E} \circ
\psi})$.

\page{37}
\note{même veine que la page 31, d'une plume plus posée : deux treillis. Dans
le premier, un cadre à main levée isole le bloc des colonnes en $\Omega$ ;
les flèches verticales de la colonne de droite sont mises entre parenthèses
sur la page. Entre les deux treillis, un bloc entièrement griffonné — on y
devine les mêmes lettres — et, aux marges, des fragments que nous
transcrivons à la suite sans pouvoir les raccorder}

\begin{tikzcd}[column sep=small, row sep=small, nodes={font=\scriptsize}]
S \arrow[d] & X' \arrow[l] \arrow[d] & F' \arrow[l] \arrow[d] &
\Omega S \arrow[l] \arrow[d] & \Omega X' \arrow[l] \arrow[d] &
\Omega F' \arrow[l] \arrow[d] & \Omega^{2} S \arrow[l] \arrow[d] \\
F \arrow[d] & G \arrow[l] \arrow[d] & \Phi \arrow[l] \arrow[d] &
\Omega F \arrow[l] \arrow[d] & \Omega G \arrow[l] \arrow[d] &
\Omega\Phi \arrow[l] \arrow[d] & \phantom{e} \\
X \arrow[d] & Z \arrow[l] \arrow[d] & G' \arrow[l] \arrow[d] &
\Omega X \arrow[l] \arrow[d] & \Omega Z \arrow[l] \arrow[d] &
\Omega G' \arrow[l] & \phantom{e} \\
S & X' \arrow[l] & F' \arrow[l] & \Omega S \arrow[l] &
\Omega X' \arrow[l] & \Omega F' \arrow[l] & \phantom{e}
\end{tikzcd}

\note{au-dessus du bloc, deux têtes de colonnes $\Omega X$ et $\Omega Z$
descendent sur $\Omega S$ et $\Omega X'$ ; les rangées se prolongent en
pointillé vers la droite}

Fragments en marge : $\Phi \to Z \to P \leftarrow F'$ (le $P$ surchargé) ;
$F \times F' \leftarrow \ill{}$ ; $F \times F' \to P \leftarrow Z$, $P \to
s$.

Le treillis en escalier des $X_{ij}$ :

\begin{tikzcd}[column sep=small, row sep=small, nodes={font=\scriptsize}]
X_{20} \arrow[d] & & & & & & & & & \\
X_{10} \arrow[d] & & & & & & & & & \\
X_{00} & X_{01} \arrow[l] & X_{02} \arrow[l] &
\Omega X_{00} \arrow[l] \arrow[d] & \Omega X_{01} \arrow[l] \arrow[d] &
\Omega X_{02} \arrow[l] \arrow[d] & \Omega^{2} X_{00} \arrow[l] \arrow[d] &
& & \\
& & & X_{20} \arrow[d] & X_{21} \arrow[l] \arrow[d] &
X_{22} \arrow[l] \arrow[d] & \Omega X_{20} \arrow[l] \arrow[d] & & & \\
& & & X_{10} \arrow[d] & X_{11} \arrow[l] \arrow[d] &
X_{12} \arrow[l] \arrow[d] & \Omega X_{10} \arrow[l] \arrow[d] & & & \\
& & & X_{00} & X_{01} \arrow[l] & X_{02} \arrow[l] &
\Omega X_{00} \arrow[l] \arrow[d] & & & \\
& & & & & & X_{20} \arrow[d] & & & \\
& & & & & & X_{10} \arrow[d] & & & \\
& & & & & & X_{00} & X_{01} \arrow[l] & X_{02} \arrow[l] &
\Omega X_{00} \arrow[l]
\end{tikzcd}

\note{deux petits arcs coiffent le $\Omega X_{00}$ et le $\Omega X_{02}$ de
la première rangée ; l'escalier dit que les suites de fibrations des lignes
et des colonnes d'un carré $X_{ij}$ s'engrènent, chaque tour de trois
rangées reprenant la précédente sous $\Omega$}

\section*{« Esquisse d'une théorie des Gr-catégories », tapuscrit annoté
(pages 39 et 40)}

\note{les pages 39 et 40 sont une copie, corrigée à l'encre, d'un tapuscrit ;
nous transcrivons le texte tel que corrigé, chaque intervention manuscrite
étant signalée ; la main des corrections n'est pas identifiée par la page.
Le texte se poursuit au-delà de la page 40, dans le lot 3}

\page{39}
\textbf{Esquisse d'une théorie des Gr-Catégories}

par Mme Hoang Xuan Sinh

(Note présentée par M. Henri Cartan)

Nous donnons un résumé de quelques résultats sur les (Gr)-catégories,
faisant l'objet d'un travail détaillé que l'auteur compte présenter
prochainement comme thèse de doctorat [ ].
\note{les crochets des renvois bibliographiques sont restés vides à la
frappe, ici et plus bas}

1. Structure des (Gr)-catégories.

Notre terminologie est celle de Saavedra [ ]. Nous nous intéressons à des
catégories $C$ munies d'une opération binaire $(X,Y) \mapsto X \otimes Y$
(foncteur de $C \times C$ dans $C$) associative et unitaire à isomorphisme
donné près (satisfaisant des conditions de compatibilité explicitées dans
loc.\ cit.), appelées aussi $\otimes$-catégories AU
(associatives-unitaires). On dit qu'une telle catégorie est une
(Gr)-catégorie si c'est un groupoïde, et si tout objet $X$ de $C$ est
« inversible », i.e.\ admet un objet « inverse » $Y = X^{-1}$ (satisfaisant
$X \otimes Y \simeq Y \otimes X \simeq 1_{C}$, où $1_{C}$ désigne l'objet
unité de $C$). Les exemples abondent :
\note{« c'est un groupoïde, et si » est un ajout à l'encre au-dessus de la
ligne}

Exemple\struck{s} 1. $X$ étant un espace topologique ponctué par $x \in X$,
on prend pour $C$ la catégorie des lacets de $X$ en $x$, avec comme
morphismes les classes d'homotopie d'homotopies entre lacets, comme
opération $\otimes$ la composition des lacets (qui n'est pas associative,
mais associative « à homotopie près »).
\note{« ponctué par $x \in X$ » remplace, à l'encre, une virgule de la
frappe}
\marginal{Cet exemple est un cas particulier du suivant, en prenant $E =
\Omega^{1}(X,x)$ (espace des lacets).}

Variante : $E$ est un espace de Hopf associatif (ou simplement
homotopiquement associatif en un sens suffisamment fort), $C$ la catégorie
dont les objets sont les points de $E$, les morphismes les classes
d'homotopies de chemins entre points de $E$, la loi $\otimes$ étant induite
par la loi de composition de $E$. (NB Lorsque $E$ admet un espace
classifiant $X$, on retrouve essentiellement l'exemple précédent).
\note{quelques mots tapés à la suite de la parenthèse sont raturés à
l'encre, illisibles}

\struck{Exemple} 2. Si $F$ est un faisceau abélien sur un espace
topologique (ou plus généralement sur un topos), la catégorie $C$ des
torseurs sous $F$, munie de la composition de Baer, est une (Gr)-catégorie
(et même une catégorie de Picard stricte, cf.\ plus bas). On peut considérer
la catégorie $C$ des Modules inversibles sur un espace (ou topos)
localement annelé $(X, O_{X})$ comme le cas particulier correspondant au cas
$F = O_{X}^{*}$.
\note{« abélien » et « localement » sont des ajouts à l'encre ; le mot
« Exemple » est rayé ici et au n° 3, ne laissant que le numéro}

\struck{Exemple} 3. Si $A$ est une catégorie, la sous-catégorie pleine $C$
de $\mathrm{Hom}(A,A)$, formée des équivalences de $A$ avec elle-même, munie
de l'opération de composition des foncteurs, est une (Gr)-catégorie. Au lieu
de prendre pour $A$ une catégorie, on peut plus généralement prendre pour
$A$ un objet d'une 2-catégorie quelconque. Si p.ex.\ $F$ est un faisceau en
groupes (pas nécessairement commutatifs) sur un espace topologique (ou un
topos \struck{$X$}),
\note{« $C$ » et « plus généralement » sont des ajouts à l'encre ; le
« $X$ » final est barré d'une grande croix}

\page{40}
prenant pour $A$ \uncertain{le} « champ » sur $X$ formé des $F$-torseurs à
droite, la (Gr)-catégorie des auto-équivalences de $A$ avec lui-même
s'interprète comme la catégorie des « bitorseurs » sous $F$, i.e.\ des
faisceaux $P$ sur lesquels $F$ opère à la fois à droite et à gauche, ces
opérations commutant et chacune d'elles faisant de $P$ un torseur (à droite
ou à gauche) sous $F$, — la composition $\otimes$ étant la composition de
Baer évidente. Lorsque $F$ est encore de la forme $O_{X}^{*}$ ($O_{X}$ étant
un faisceau d'anneaux, qu'on ne suppose plus nécessairement commutatif) ces
bitorseurs s'interprètent aussi en termes de bi-Modules « inversibles » sous
$O_{X}$.
\note{la frappe porte « la "champ" » ; le $P$ après « faisceaux » est un
ajout à l'encre}

\textbf{Structure.} Soit $C$ une (Gr)-catégorie, on lui associe

a) le groupe $\pi_{0}(C)$ des classes d'isomorphisme d'objets de $C$,

b) le groupe $\pi_{1}(C)$ des automorphismes de $1_{C}$ (objet unité de
$C$)

c) une action de $\pi_{0}(C)$ sur $\pi_{1}(C)$, en associant à tout objet
$X$ de $C$ l'automorphisme $\rho(X)$ de $1_{C}$ déduit des deux
isomorphismes
\[
\mathrm{Aut}(1_{C}) \rightrightarrows \mathrm{Aut}(X)
\]
donnés par $u \mapsto u \otimes \mathrm{id}_{X}$ et $u \mapsto
\mathrm{id}_{X} \otimes u$. On prouve que $\pi_{1}(C)$ est un groupe
\emph{commutatif} et que l'on obtient bien par c) une structure de
$\pi_{0}(C)$-module sur celui-ci. Ceci posé, si on choisit pour tout $a \in
\pi_{0}(C)$ un représentant $L_{a} \in \mathrm{Ob}\,C$ de $a$, et pour deux
\struck{objets} $a, b$ un isomorphisme
\[
\phi_{a,b} : L_{a} \otimes L_{b} \simeq L_{ab},
\]
alors pour trois éléments $a, b, c \in \pi_{0}(C)$, l'isomorphisme
d'associativité
\[
(L_{a} \otimes L_{b}) \otimes L_{c} \simeq L_{a} \otimes (L_{b} \otimes
L_{c}),
\]
compte tenu des isomorphismes $\phi_{a,b}$, $\phi_{ab,c}$, $\phi_{b,c}$ et
$\phi_{a,bc}$, peut s'interpréter comme un isomorphisme
\[
L_{abc} \simeq L_{abc},
\]
ou encore comme la tensorisation à gauche avec un élément bien déterminé
\[
f(a,b,c) \in \pi_{1}(C).
\]
Les $\phi_{a,b}$ étant choisis, on voit donc que la donnée d'un isomorphisme
d'associativité fonctoriel $(L \otimes L') \otimes L'' \simeq L \otimes (L'
\otimes L'')$ équivaut à la donnée de l'application
\[
f : \pi_{0} \times \pi_{0} \times \pi_{0} \longrightarrow \pi_{1}.
\]
On vérifie alors que l'axiome standard d'autocompatibilité d'une donnée
d'associativité (axiome des pentagones) s'exprime précisément par la
condition que $f$ soit un \emph{3-cocycle} du groupe $\pi_{0}$ à valeurs
dans le groupe $\pi_{1}$. D'autre part, l'indétermination dans le choix du
système d'isomorphismes $\phi_{a,b}$ est précisément donnée par une
2-cochaine arbitraire, et on voit que si on change $\phi$ par une 2-cochaine
$g$, $f$ est changé en $f + dg$ — donc l'ensemble des $f$ correspondants à
des choix différents de $\phi$ est exactement une classe de 3-cohomologie
\note{« (axiome des pentagones) » est un ajout à l'encre ; la page
s'interrompt sans ponctuation — la suite du tapuscrit est au lot 3}

\end{document}
